% Requires amsmath, amssymb, amsthm; no private macros.
\section{F-23 over every algebraically closed field}
\label{sec:fields}

Throughout this section $k$ is algebraically closed, with arbitrary
characteristic. Border rank means CP border rank in the displayed
tensor factors, defined by Zariski closure of sums of pure tensors.
The integral convention below is essential: the normalized
characteristic-zero coefficient tensor cannot simply be reduced
modulo a prime dividing its normalization denominator.

\paragraph{Finding FL-01 (the integral models).}
Let $A_n=\mathbb Z^n$ have basis $e_1,\ldots,e_n$, and let
$V_n=\mathbb Z^{n^2}$ have basis $x_{ij}$.
Define
\begin{align*}
 D_n&=\sum_{\sigma\in\mathfrak S_n}\operatorname{sgn}(\sigma)
             e_{\sigma(1)}\otimes\cdots\otimes e_{\sigma(n)},&
 P_n&=\sum_{\sigma\in\mathfrak S_n}
             e_{\sigma(1)}\otimes\cdots\otimes e_{\sigma(n)},\\
 S_{\det,n}&=\sum_{\sigma,\tau\in\mathfrak S_n}
       \operatorname{sgn}(\sigma)\operatorname{sgn}(\tau)
       x_{\sigma(1),\tau(1)}\otimes\cdots\otimes x_{\sigma(n),\tau(n)},\\
 S_{\mathrm{per},n}&=\sum_{\sigma,\tau\in\mathfrak S_n}
       x_{\sigma(1),\tau(1)}\otimes\cdots\otimes x_{\sigma(n),\tau(n)}.
\end{align*}
The latter two are invariant under permuting the tensor slots.
For free modules they therefore belong to the integral divided-power
model $\Gamma^n(V_n)\simeq(V_n^{\otimes n})^{\mathfrak S_n}$.
In characteristic zero they are $n!$ times the normalized polynomial
coefficient tensors of predecessor GCT\_GINEQ F-23. Their images in
the ordinary symmetric quotient are $n!\det_n$ and $n!\mathrm{per}_n$.

\begin{proof}
An ordered sequence of matrix positions with every row and every
column used once uniquely determines its row permutation $\sigma$
and column permutation $\tau$. Thus each tensor-basis word occurs
once, with coefficient $\pm1$ or $1$. Permuting all slots composes
both permutations by the same permutation, so their sign product
is unchanged. A basis of the invariant tensors consists of the sums
over distinct reorderings of each monomial; this is the usual
divided-power basis. In the ordinary symmetric quotient, each
squarefree matching monomial has $n!$ such reorderings. The claimed
formulas follow. In characteristic two signs coalesce, so the
determinant and permanent models agree. No integral coefficient
vanishes by normalization: each surviving ordered word has a unit
coefficient. If $p\mid n!$, the symmetric quotient of either
invariant tensor is zero while the invariant tensor is nonzero.
\end{proof}

These definitions fix what is meant by the polynomial families in
an all-field statement. They do not identify the ordinary symmetric
quotient with invariant tensors by division by $n!$.

\paragraph{Finding FL-02 (proper-incidence border substitution).}
Let $T\in A\otimes V_1\otimes\cdots\otimes V_b$, let
$W=V_1\otimes\cdots\otimes V_b$, and let $U\subset W$ be the
image of the $A$-flattening. If $q=\dim U\geq1$, then
\begin{equation}\label{eq:fields-substitution}
 \underline{\operatorname R}_k(T)
 \geq q-1+\min_{0\ne u\in U}\underline{\operatorname R}_k(u).
\end{equation}
This is the needed all-field, arbitrary-residual-order form of
border substitution. Landsberg--Micha\l ek
\cite[Proposition~2.3, p.~3]{LM2018} already state the three-factor
principle for every retained dimension, not just dimension one.
The all-field, arbitrary-residual-order form follows from R-01.

\begin{proof}
Apply R-01 with retained dimension $j=1$. A line
$L=\langle u\rangle$ has inclusion tensor $a\otimes u$,
whose border rank equals that of $u$: adjoining a fixed nonzero
factor and contracting against its dual are inverse linear maps
on this tensor. Consequently the one-dimensional Grassmannian
minimum in R-01 is precisely the minimum in
\eqref{eq:fields-substitution}.
\end{proof}

\paragraph{Finding FL-03 (exact residual minima and recursion).}
Let $T_n$ denote any one of the four integral families in FL-01.
For $2\leq b<n$, group the first $n-b$ slots and denote the
remaining flattening image by $U_{n,b}$. Put $s=1$ for the row
families and $s=2$ for the polynomial divided-power families.
Then
\begin{align}
 \min_{0\ne u\in U_{n,b}}\underline{\operatorname R}_k(u)
     &=\underline{\operatorname R}_k(T_b),\label{eq:fields-minimum}\\
 \underline{\operatorname R}_k(T_n)
     &\geq\binom nb^s+\underline{\operatorname R}_k(T_b)-1.
       \label{eq:fields-recursion}
\end{align}

\begin{proof}
In a row family, contract the first $n-b$ slots against an ordered
tuple of coordinate covectors. If an index repeats, the result is
zero. Otherwise the result is, up to a sign, the embedded copy
$Q_I$ of $T_b$ on the complementary $b$-subset $I$. Every $I$
occurs. Distinct $Q_I$ have disjoint tensor-basis supports, each
with a unit coefficient. They therefore form a free integral basis
of the image, remaining independent after every scalar extension.
Its dimension is $\binom nb$.

For the polynomial families, a nonzero contraction is exactly a
partial row-and-column matching. Its residual is, up to a sign,
the coefficient tensor $Q_{I,J}$ of the complementary minor or
permanent. Every pair of $b$-subsets $I,J$ occurs. The residual
supports are disjoint and have unit coefficients, giving image
dimension $\binom nb^2$ over every field. This is the tensor version
of the complementary-minor calculation in GCT\_GINEQ F-22.
In particular, no factor $b!$ or $(n-b)!$ occurs here.

Choose a nonzero coefficient $c$ of $u$ in this basis. In a row
family, project every residual factor onto the coordinates $I$
of that coefficient. Each other $b$-subset has a required index
outside $I$ and is killed, whereas $Q_I$ maps to $T_b$.
In a polynomial family, project onto the matrix-entry block
$I\times J$ in every factor. Each other pair of $b$-subsets has
a required row or column outside this block, and is killed.
The image is $cT_b$.

Products of linear maps preserve pure tensors or send them to
zero; continuity gives the same border-rank monotonicity.
Consequently every nonzero $u$ has border rank at least that of
$T_b$. Conversely, a single $Q_I$ or $Q_{I,J}$ is an embedding of
$T_b$ with the displayed projection as a left inverse, so its
border rank is exactly that of $T_b$. This proves
\eqref{eq:fields-minimum}. Grouping the first $n-b$ factors
cannot increase CP border rank. Apply FL-02 to that grouped
tensor and the image dimension just computed to obtain
\eqref{eq:fields-recursion} for the original tensor.
\end{proof}

\paragraph{Finding FL-04 (unit-minor seeds).}
Over every algebraically closed field,
\[
 \underline{\operatorname R}_k(D_2)=
 \underline{\operatorname R}_k(P_2)=2,\qquad
 \underline{\operatorname R}_k(S_{\det,2})=
 \underline{\operatorname R}_k(S_{\mathrm{per},2})=4.
\]
Hence, for $n\geq2$,
\[
 \underline{\operatorname R}_k(D_n),
 \underline{\operatorname R}_k(P_n)\geq\binom n2+1,
 \qquad
 \underline{\operatorname R}_k(S_{\det,n}),
 \underline{\operatorname R}_k(S_{\mathrm{per},n})
       \geq\binom n2^2+3.
\]

\begin{proof}
The two row matrices are
$\left(\begin{smallmatrix}0&1\\-1&0\end{smallmatrix}\right)$
and $\left(\begin{smallmatrix}0&1\\1&0\end{smallmatrix}\right)$.
The polynomial matrices have two off-diagonal coordinate pairs,
$(x_{11},x_{22})$ and $(x_{12},x_{21})$, with unit entries.
They are signed permutation matrices of size four. Their ranks
are respectively two and four over every field. Matrix rank
equals matrix border rank. For $n>2$ use FL-03 with $b=2$;
for $n=2$ use the seeds themselves.
\end{proof}

Thus the predecessor's bound $39$ at $n=4$ has an all-field
proof in the explicit integral model. Stronger multidimensional
certificates elsewhere in the paper improve this seed recursion
over every algebraically closed field in the integral convention.

\paragraph{Finding FL-05 (ordinary polynomial powers in small characteristic).}
If $\operatorname{char}k=p\leq n$, neither ordinary polynomial
$\det_n$ nor $\mathrm{per}_n$ lies in any secant of the image
$\ell\mapsto\ell^n$ in $\operatorname{Sym}^n(V_n)$.
Their border rank relative to this ordinary-power image is
therefore infinite. This statement concerns the ordinary symmetric
quotient, not the complete-linear-system Veronese in divided powers.

\begin{proof}
The coefficient of a squarefree degree-$n$ monomial in $\ell^n$
is $n!$ times a product of coefficients of $\ell$. It is zero
when $p\leq n$. All sums of ordinary powers, and all their limits,
therefore lie in the closed linear subspace on which every such
coefficient vanishes. Determinant and permanent each have nonzero
squarefree matching coefficients and lie outside that subspace.
\end{proof}

When $n!$ is invertible, the usual symmetrization identification
is valid. In that range, every CP border-rank lower bound here
also bounds ordinary polynomial symmetric border rank. In small
characteristic the finite meaningful assertion uses FL-01's
divided-power model and its complete Veronese; FL-05 explains why
the alternative ordinary-power convention has a different answer.
